\section*{Definitions and notation (author source p. 13:3)}
Given a Boolean matrix $M\in\{0,1\}^{m\times n}$, let us define its 1-partition number, $\chi_1(M)$, as the smallest $s$ such that $M$ can be written as
\[M=LR,\quad\text{with }L\in\{0,1\}^{m\times s},\ R\in\{0,1\}^{s\times m},\]
where the operations are over $\RR$. The definition emphasizes the analogy with $\rk_+$, and $\chi_1$ is also sometimes referred to as binary rank. On the other hand, the phrase ``partition number'' comes from communication complexity. The name is justified: it is easy to see that $\chi_1(M)$ equals the smallest $s$ such that the 1-entries of $M$ can be partitioned into $s$ 1-monochromatic rectangles (i.e., rank-one Boolean matrices). Finally, when $M$ is viewed as the adjacency matrix of a bipartite graph, $\chi_1(M)$ also appears under the name biclique partition number [3].
\paragraph{Notation.} All logarithms are in base 2 and $[n]:=\{1,\ldots,n\}$.
\section*{4 1-Partition number (author source pp. 13:4--13:5)}
In this section, we prove Theorem 2.
Let $M$ be an $m\times n$ matrix with rows indexed by $[n]=\{1,\ldots,n\}$. Given $A\subseteq[n]$, $M_A$ denotes the submatrix obtained by removing the rows outside of $A$ from $M$. Observe that\footnote{If $A_1,A_2$ are disjoint, this is quite obvious. Otherwise consider $A_1,A_2\setminus A_1$.}
\[\chi_1(M_{A_1\cup A_2})\le\chi_1(M_{A_1})+\chi_1(M_{A_2}),\tag{1}\]
and so $\chi_1(M_A)$ can be viewed as a subadditive measure on $[n]$ whenever $M$ is fixed.
If a matrix $M$ has rank $r$, its rows are a linear combination of a subset of $r$ rows of $M$. This means that every column of $M$ is determined by a fixed subset of $r$ coordinates. If $M$ is Boolean, this leads to the following useful fact:
if $M$ has distinct columns then $n\le 2^{\rk(M)}$ (similarly for rows).
\begin{originalresult}{Lemma 5}
Let $M$ be an $m\times n$ Boolean matrix of rank $r$. Given $s\in[n]$, let $\mathcal A$ be the collection of maximal subsets $A\subseteq[n]$ with $\chi_1(M_A)\le s$ (i.e., $\chi_1(M_A)\le s$ and $\chi_1(M_{A'})>s$ for every $A'\supsetneq A$). Then $|\mathcal A|\le 2^{(r+s)^2}$.
\end{originalresult}
\begin{proof}
Let $v_1,\ldots,v_n\in\RR^m$ be the columns of $M$. Given $L\in\{0,1\}^{m\times s}$, let
\[L^*:=\{i\in[n]:\exists y\in\{0,1\}^s\ v_i=Ly\}.\]
Let $\mathcal L:=\{L^*:L\in\{0,1\}^{m\times s}\}$.
We claim that $\mathcal A\subseteq\mathcal L$. If $\chi_1(M_A)\le s$, we can write $M_A=LR$ with $L\in\{0,1\}^{m\times s}$ and $R\in\{0,1\}^{s\times|A|}$. This means that every $v_i$, $i\in A$, is a Boolean linear combination of the columns of $L$ and $A\subseteq L^*$. Furthermore, if $A$ is maximal, we must have $A=L^*$.
We now want to estimate the size of $\mathcal L$. The set $L^*$ consists of indices $i\in[n]$ so that there exists $x\in\RR^n,y\in\RR^s$ satisfying
\[Mx-Ly=0\tag{2}\]
such that $y\in\{0,1\}^s$ and $x$ is the $i$-th unit vector. Since $M$ has rank $r$ and $L$ has rank at most $s$, the system (2) is equivalent to a subsystem of $t:=\min((s+r),m)$ equations. These correspond to rows of the matrix $(M,L)$. Hence, in order to determine $L^*$, it is sufficient to specify a $t$-element subset of $[m]$ together with the $t\times s$ submatrix of $L$. This gives the estimate
\[|\mathcal L|\le\binom mt2^{ts}\le2^{t(s+\log m)}.\]
Finally, we can assume that $M$ has distinct rows and so $\log m\le r$, obtaining the bound $2^{(r+s)^2}$.
\end{proof}
